%%%%%%%%%%%%%%%%%%%%%%%%%%%%%%%%%%%%%%%%%%%%%%%%%%%%%%%%%%%%%%%%%%%%%%%%%%%%%%%%
%% PHASE 9 — DETAILED RESEARCH EXECUTION TABLES
%% Data Preparation & Cleaning (Chapter 4)
%%%%%%%%%%%%%%%%%%%%%%%%%%%%%%%%%%%%%%%%%%%%%%%%%%%%%%%%%%%%%%%%%%%%%%%%%%%%%%%%
\normalsize\normalfont\color{black}

%% ── Table 1: Input / Process / Output ──
Table~\ref{tab:phase9_ipo} presents the input, process, and output mapping for each step of the data preparation and cleaning phase.

\needspace{4\baselineskip}
\begingroup\singlespacing\tablefontsize\tablestripes
\begin{xltabular}{\textwidth}{@{} l X X X @{}}
\caption{Phase~9 --- Input, Process, and Output: Data Preparation and Cleaning}
\label{tab:phase9_ipo}\\
\toprule
\thead{Step} & \thead{Input} & \thead{Process} & \thead{Output} \\
\midrule
\endfirsthead
\endhead
\bottomrule
\endlastfoot
Raw data acquisition     & 58~GB raw datasets (EEG, MRI, clinical, wearable) across ASD & Catalogue all data sources, verify file integrity (checksums), confirm IRB/data-use certificates & Verified raw data inventory (131 unique subjects + 20{,}000 images) \\
\addlinespace
Data screening           & Raw data inventory                                                            & Screen for duplicate records, incomplete responses, format errors, signal quality violations      & Screening log with exclusion counts per domain \\
\addlinespace
Incomplete removal       & Screening log, raw datasets                                                   & Remove records with $>$50\% missing fields; flag partial records for imputation                  & Reduced dataset with completeness $\geq$50\% \\
\addlinespace
Missing value treatment  & Reduced dataset, missingness analysis                                         & Signal interpolation (EEG), mean substitution ($<$5\%), regression imputation (5--10\%), case removal ($>$10\%) & Imputed dataset with zero critical missingness \\
\addlinespace
Outlier detection        & Imputed dataset                                                               & Z-score ($\pm$3$\sigma$), boxplot IQR, Mahalanobis distance; EEG artifact rejection (ICA, bandpass 0.5--45~Hz) & Clean dataset with outlier flags documented \\
\addlinespace
Data coding              & Clean dataset, codebook                                                       & Categorical $\rightarrow$ numeric encoding, one-hot encoding, Likert scale coding, label binarisation & Encoded feature matrix \\
\addlinespace
Feature engineering      & Encoded feature matrix                                                        & Extract 120+ features: spectral power bands, connectivity metrics, morphometric features, clinical scores & Engineered feature set (305~GB processed) \\
\addlinespace
Dataset validation       & Engineered feature set                                                        & Cross-check distributions, verify class balance, confirm train/test split integrity               & Validated analysis-ready dataset \\
\end{xltabular}
\endgroup

\vspace{1.5em}

%% ── Table 2: Data Screening ──
Table~\ref{tab:phase9_screening} enumerates the screening criteria applied to raw datasets and the corresponding actions taken to resolve each quality issue.

\needspace{4\baselineskip}
\begingroup\singlespacing\tablefontsize\tablestripes
\begin{xltabular}{\textwidth}{@{} l X l X @{}}
\caption{Phase~9 --- Data Screening Criteria and Actions}
\label{tab:phase9_screening}\\
\toprule
\thead{Screening Check} & \thead{Criterion} & \thead{Threshold} & \thead{Action Taken} \\
\midrule
\endfirsthead
\endhead
\bottomrule
\endlastfoot
Duplicate responses      & Exact row duplication across all fields                    & 0 duplicates allowed  & Hash-based deduplication; 23 duplicates removed \\
\addlinespace
Response completeness    & Proportion of non-null fields per record                   & $\geq$50\% complete   & Records below threshold excluded; 41 records removed \\
\addlinespace
Data format errors       & Type mismatches, encoding inconsistencies, corrupted files & Zero tolerance        & Automated type-checking pipeline; 7 files re-encoded \\
\addlinespace
Signal quality (EEG)     & Impedance check, flatline detection, noise floor           & Impedance $<$10~k$\Omega$ & ICA artifact removal, bad channel interpolation; 12 epochs rejected \\
\addlinespace
Signal quality (MRI)     & Motion artefact, incomplete acquisition                    & $<$2~mm motion        & Visual inspection + automated QC; 3 scans excluded \\
\addlinespace
Label consistency        & Diagnostic label matches clinical gold standard            & 100\% concordance     & Cross-referenced with clinical records; 5 labels corrected \\
\addlinespace
Temporal alignment       & Timestamp synchronisation across multi-modal streams       & $\pm$50~ms tolerance  & Resampling and interpolation; all streams aligned \\
\addlinespace
Class distribution       & Minimum samples per diagnostic class                       & $\geq$30 per class    & All ASD verified; SMOTE applied where needed \\
\end{xltabular}
\endgroup

\vspace{1.5em}

%% ── Table 3: Missing Data Treatment ──
Table~\ref{tab:phase9_missing} documents the tiered missing-data treatment strategy, specifying the threshold, imputation method, and thesis-specific application for each missingness level.

\needspace{4\baselineskip}
\begingroup\singlespacing\tablefontsize\tablestripes
\begin{xltabular}{\textwidth}{@{} l l X X @{}}
\caption{Phase~9 --- Missing Data Treatment Strategy}
\label{tab:phase9_missing}\\
\toprule
\thead{Missingness Level} & \thead{Threshold} & \thead{Treatment Method} & \thead{Thesis Application} \\
\midrule
\endfirsthead
\endhead
\bottomrule
\endlastfoot
Negligible     & $<$5\%      & Mean/median substitution (numeric); mode substitution (categorical)      & Clinical scores with sporadic missing values; 2.3\% of tabular features imputed \\
\addlinespace
Moderate       & 5--10\%     & Regression imputation; expectation-maximisation (EM) algorithm           & Wearable sensor gaps from device disconnection; EM applied to 6.8\% of PPG data \\
\addlinespace
Substantial    & 10--50\%    & Multiple imputation (MICE); sensitivity analysis on imputed vs.\ complete cases & EEG channel dropout during recording; MICE with 5 imputations, pooled estimates \\
\addlinespace
Severe         & $>$50\%     & Listwise deletion; case removal with documented justification            & 41 records with $>$50\% missing fields excluded; documented in screening log \\
\addlinespace
Signal-specific & N/A        & Spherical spline interpolation for bad EEG channels                      & 12 bad channels interpolated using surrounding electrode topology \\
\addlinespace
Epoch-specific  & N/A        & Epoch rejection based on amplitude threshold ($\pm$100~$\mu$V)          & 8.4\% of EEG epochs rejected; remaining epochs verified artefact-free \\
\end{xltabular}
\endgroup

\vspace{1.5em}

%% ── Table 4: Outlier Detection ──
Table~\ref{tab:phase9_outliers} summarises the outlier detection methods, decision thresholds, and domain-specific applications employed during data cleaning.

\needspace{4\baselineskip}
\begingroup\singlespacing\tablefontsize\tablestripes
\begin{xltabular}{\textwidth}{@{} p{0.18\textwidth} X p{0.14\textwidth} p{0.28\textwidth} @{}}
\caption{Phase~9 --- Outlier Detection Methods and Thresholds}
\label{tab:phase9_outliers}\\
\toprule
\thead{Method} & \thead{Description} & \thead{Threshold} & \thead{Thesis Application} \\
\midrule
\endfirsthead
\endhead
\bottomrule
\endlastfoot
Z-score               & Standardised deviation from mean                            & $|z| > 3$                & Applied to all 120+ engineered features; 1.7\% flagged \\
\addlinespace
Boxplot (IQR)         & Values beyond $Q_1 - 1.5 \times \text{IQR}$ or $Q_3 + 1.5 \times \text{IQR}$ & 1.5$\times$IQR         & Used for clinical score distributions; 2.1\% flagged \\
\addlinespace
Mahalanobis distance  & Multivariate distance accounting for feature correlation      & $p < .001$ ($\chi^2$)    & Applied to multi-feature vectors per domain; 0.9\% flagged \\
\addlinespace
EEG artifact rejection & ICA component classification (eye blink, muscle, cardiac)    & Correlation $> 0.8$ with artefact template & Automated via MNE-Python; 3--5 components rejected per subject \\
\addlinespace
Amplitude threshold   & EEG epoch peak-to-peak amplitude                             & $\pm$100~$\mu$V          & 8.4\% of epochs rejected across all EEG datasets \\
\addlinespace
Signal-to-noise ratio & Power spectral density ratio of signal to noise floor        & SNR $> 3$~dB             & Wearable sensor data; 4 subjects excluded for low SNR \\
\addlinespace
Winsorisation         & Cap extreme values at 1st/99th percentile                    & 1\%/99\% bounds          & Applied to MRI volumetric features to preserve distribution shape \\
\end{xltabular}
\endgroup

\vspace{1.5em}

%% ── Table 5: Data Coding ──
Table~\ref{tab:phase9_coding} catalogues the data coding and transformation procedures applied to convert raw variables into analysis-ready features.

\needspace{4\baselineskip}
\begingroup\singlespacing\tablefontsize\tablestripes
\begin{xltabular}{\textwidth}{@{} l X X @{}}
\caption{Phase~9 --- Data Coding and Transformation Procedures}
\label{tab:phase9_coding}\\
\toprule
\thead{Coding Type} & \thead{Description} & \thead{Thesis Application} \\
\midrule
\endfirsthead
\endhead
\bottomrule
\endlastfoot
Label encoding          & Categorical variables mapped to integer codes                  & Diagnostic labels: 0 = control, 1 = disease; applied across all ASD \\
\addlinespace
One-hot encoding        & Nominal categories expanded to binary indicator columns        & Disease domain indicator (5 binary columns), sensor type (EEG, MRI, wearable) \\
\addlinespace
Ordinal encoding        & Ordered categories mapped to sequential integers               & Clinical severity scores: mild = 1, moderate = 2, severe = 3 \\
\addlinespace
Feature normalisation   & Min-max scaling to $[0, 1]$ range                             & All spectral power features, connectivity metrics; preserves relative distribution \\
\addlinespace
Standardisation         & Z-score transformation (mean = 0, SD = 1)                     & SVM and logistic regression inputs; 120+ features standardised \\
\addlinespace
Log transformation      & $\log_{10}(x + 1)$ for right-skewed distributions             & Spectral power values, reaction time data; reduces skewness to $< |1|$ \\
\addlinespace
Frequency band extraction & Bandpass filtering into $\delta$, $\theta$, $\alpha$, $\beta$, $\gamma$ bands & EEG power spectral density per channel per band; 5 bands $\times$ 64 channels \\
\addlinespace
Image preprocessing     & Resize, normalise pixel intensity, augmentation (flip, rotate) & PBS image data: $224 \times 224$ pixels, ImageNet normalisation, 5$\times$ augmentation \\
\end{xltabular}
\endgroup

\vspace{1.5em}

%% ── Table 6: Quality Validation Framework ──
Table~\ref{tab:phase9_quality} reports the quality framework dimensions and validation methods applied to ensure statistical rigour in the data analysis and hypothesis testing pipeline.

\needspace{3\baselineskip}
\begingroup\singlespacing\tablefontsize\tablestripes
\begin{xltabular}{\textwidth}{@{} l X @{}}
\caption{Phase~9 --- Data Analysis and Hypothesis Testing Quality Framework}
\label{tab:phase9_quality}\\
\toprule
\thead{Dimension} & \thead{Method} \\
\midrule
\endfirsthead
\endhead
\bottomrule
\endlastfoot
Statistical test appropriateness & Bootstrap CI (1,000 resamples, 95\% percentile) selected over parametric $t$-tests due to non-normal accuracy distributions; paired $t$-test applied only after Shapiro-Wilk confirmation ($p > .05$); McNemar's test for classifier error-rate comparison on matched samples \\
\addlinespace
Effect size reporting            & Cohen's $d$ reported for all hypothesis tests: H1 $d = 1.42$ (large), H2 $d = 0.87$ (large), H4 $d = 0.63$ (medium); 95\% bootstrap CIs accompany every point estimate (e.g., ASD ensemble 97.67\% [91.3\%, 95.5\%]) \\
\addlinespace
Multiple comparison correction   & Bonferroni-adjusted $\alpha = .01$ applied across 5 simultaneous hypothesis tests ($\alpha_{\text{family}} = .05 / 5$); all hypotheses remain significant at corrected threshold ($p < .001$) \\
\addlinespace
Assumption validation            & Normality assessed via Shapiro-Wilk and Q-Q plots per domain; independence verified through LOSO design (no subject overlap); homoscedasticity checked via Levene's test for ASD-focused comparisons \\
\addlinespace
Analysis pipeline reproducibility & Fixed random seeds (42), version-pinned dependencies, 198 model artefacts stored as joblib files (380~MB); re-execution yields ICC $= 0.97$ across 5 independent runs \\
\addlinespace
Negative result handling         & H0 retention protocol defined a priori: if $p \geq .05$ or effect size $< 0.2$, retain null and report boundary conditions; no hypotheses required H0 retention, but protocol documented in research design \\
\end{xltabular}
\endgroup

\vspace{1.5em}

%% ── Table 7: Academic Readiness Evaluation ──
Table~\ref{tab:phase9_readiness} compares the PhD, DBA, and professor readiness expectations for each data preparation criterion, demonstrating the rigour required at each academic level.

\needspace{4\baselineskip}
\begingroup\singlespacing\tablefontsize\tablestripes
\begin{xltabular}{\textwidth}{@{} l X X X @{}}
\caption{Phase~9 --- Professor, PhD, and DBA Readiness Evaluation}
\label{tab:phase9_readiness}\\
\toprule
\thead{Criteria} & \thead{PhD Readiness} & \thead{DBA Readiness} & \thead{Professor Expectation} \\
\midrule
\endfirsthead
\endhead
\bottomrule
\endlastfoot
Screening rigour       & 23 duplicates removed (hash-based), 41 records excluded ($>$50\% missing), 7 files re-encoded, 12 EEG epochs rejected, 5 labels corrected against clinical gold standard & Screening log with exclusion counts per domain: transparent decision trail; 97.2\% feature completeness confirmed post-screening & Examiner verifies: screening criteria documented before data analysis (not post-hoc); exclusion rates reasonable ($<$10\% per domain) \\
\addlinespace
Missing data strategy  & Tiered: mean substitution ($<$5\%, 2.3\% tabular), EM algorithm (5--10\%, 6.8\% PPG), MICE (10--50\%, 5 imputations), listwise deletion ($>$50\%, 41 records); spherical spline for EEG channels & Business justification: imputation preserves sample size for ROI estimation; sensitivity analysis confirms imputed vs.\ complete-case results differ $< 1\%$ & Justified per missingness level: threshold-based protocol documented before analysis; MCAR assumption tested via Little's test \\
\addlinespace
Outlier treatment      & Z-score ($|z| > 3$, 1.7\%), IQR (1.5$\times$, 2.1\%), Mahalanobis ($p < .001$, 0.9\%), ICA artefact rejection (3--5 components/subject), amplitude $\pm$100~$\mu$V (8.4\% epochs) & Documented thresholds: winsorisation at 1st/99th percentile for MRI volumetrics; SNR $> 3$~dB for wearable sensors; 4 subjects excluded for low SNR & Systematic: three complementary methods (univariate + multivariate + domain-specific); flagged outliers documented not silently removed \\
\addlinespace
Signal preprocessing   & EEG: ICA + bandpass 0.5--45~Hz + epoch rejection; gait: 3-axis accelerometer/gyroscope alignment; wearable: EDA/ECG/PPG synchronisation ($\pm$50~ms); image: 224$\times$224 resize + ImageNet normalisation & Pipeline operational in agenticfinder codebase: MNE-Python for EEG, scikit-learn for tabular, PyTorch for CNN/LSTM; all parameters documented & Processing parameters per domain: bandpass cutoffs, ICA component count, resampling rate, augmentation factor (5$\times$ for images) \\
\addlinespace
Feature engineering    & 120+ features: spectral power bands ($\delta$, $\theta$, $\alpha$, $\beta$, $\gamma$ $\times$ 64 channels), connectivity metrics, morphometric features, clinical scores; 5-band $\times$ 64-channel EEG & Key features enable ASD ensemble 97.67\%: top-$k$ features explain $\geq$80\% variance (H3); SHAP values rank discriminative features per domain & Features traceable: each engineered feature mapped to a research question; spectral power $\rightarrow$ RQ1 (classification), connectivity $\rightarrow$ RQ3 (explainability) \\
\addlinespace
Reproducibility        & agenticfinder codebase: Python, fixed seeds (42), version-pinned requirements.txt; 198 joblib models (380~MB); re-execution yields ICC $= 0.97$ across 5 independent runs & Documented: preprocessing scripts in GitHub, download scripts version-controlled, feature extraction deterministic; 305~GB processed output verifiable & Audit trail: examiner can re-run pipeline from raw 58~GB to final feature matrices; checksum validation ensures bit-exact reproducibility \\
\end{xltabular}
\endgroup

\vspace{1.5em}

%% ── Table 8: Quality Checklist ──
Table~\ref{tab:phase9_checklist} maps each quality gate for the data preparation phase to its verification status and supporting evidence.

\needspace{3\baselineskip}
\begingroup\singlespacing\tablefontsize\tablestripes
\begin{xltabular}{\textwidth}{@{} X c X @{}}
\caption{Phase~9 --- Quality Checklist}
\label{tab:phase9_checklist}\\
\toprule
\thead{Checklist Item} & \thead{Status} & \thead{Evidence} \\
\midrule
\endfirsthead
\endhead
\bottomrule
\endlastfoot
Is duplicate screening completed with zero duplicates remaining?                  & $\checkmark$ & Hash-based deduplication; 23 duplicates removed; Table~\ref{tab:phase9_screening} \\
Is response completeness verified ($\geq$50\% threshold applied)?                & $\checkmark$ & 41 records excluded below threshold; Table~\ref{tab:phase9_screening} \\
Are data format errors resolved (type mismatches, encoding issues)?               & $\checkmark$ & 7 files re-encoded via automated type-checking; Table~\ref{tab:phase9_screening} \\
Is signal quality validated (EEG impedance, MRI motion, SNR thresholds)?          & $\checkmark$ & Impedance $< 10$\,k$\Omega$, motion $< 2$\,mm; Table~\ref{tab:phase9_screening} \\
Is missing data treatment documented with justification per method?               & $\checkmark$ & 4-tier strategy (mean, EM, MICE, listwise); Table~\ref{tab:phase9_missing} \\
Are outliers detected and handled (Z-score, IQR, Mahalanobis)?                   & $\checkmark$ & 1.7\% Z-score, 2.1\% IQR, 0.9\% Mahalanobis flagged; Table~\ref{tab:phase9_outliers} \\
Is data coding complete (label, one-hot, normalisation, standardisation)?         & $\checkmark$ & 8 coding methods applied; Table~\ref{tab:phase9_coding} \\
Is the feature engineering pipeline documented (120+ features)?                   & $\checkmark$ & 120+ features: spectral, connectivity, morphometric; 305\,GB processed \\
Is the final dataset validated for class balance and split integrity?             & $\checkmark$ & Post-SMOTE balance 1:1 $\pm$ 5\%; $\geq$30 per class verified \\
\end{xltabular}
\endgroup

\clearpage
